% =====================================================================
%  results_ktarget_burnup.tex
%  Version 2, 3 September 2026. FILLED with the measured results of
%  kt_burnup/summary.json (10 designs, 3 seeds, 90 core and assembly
%  solves, 8.43 h of solve wall time).
%
%  Paste into results_c8_v2.tex between the AUTHOR note that closes
%  \subsection{The two-bank subset} (around line 438) and
%  \subsection{Cycle length against the reported refuelling interval}
%  (around line 441).
%
%  Labels defined: sec:res-kt-burnup, tab:kt-burnup, eq:kt-lf,
%                  eq:kt-defpd, eq:kt-comp
%  Labels used:    sec:meth-ktarget, sec:meth-lax, sec:res-corevalidation,
%                  sec:res-c8-front, sec:res-c8-limits, tab:c8-front,
%                  tab:c8-two, eq:ktarget
%
%  Requires: booktabs, siunitx (both already in Thesis_template.tex).
%  Two AUTHOR notes remain, marked %% AUTHOR.
% =====================================================================

\subsection{Burnup and composition dependence of the leakage
correction}
\label{sec:res-kt-burnup}

The leakage-corrected target of Equation~\eqref{eq:ktarget} rests on
two assumptions that the closure check of
Section~\ref{sec:res-corevalidation} does not reach. The leakage
factor is evaluated once, at beginning of life, on a single reference
composition of \num{4.55} and \SI{4.05}{\percent} by weight without
gadolinia, and it is then held constant over the whole cycle and
applied to every design in the space. Both assumptions are tested
here, on the eight front designs of Table~\ref{tab:c8-front} and on
the two-bank representatives of Table~\ref{tab:c8-two}.

For each design the reflective assembly is depleted once, in a single
shot with the campaign schedule, up to two steps beyond its end of
cycle. At beginning of life and at the two depletion steps that
bracket the end of cycle, the depleted composition is loaded uniformly
into all \num{32} assemblies of the core model of
Section~\ref{sec:meth-ktarget} and into the reflective assembly, and
both eigenvalues are solved with three independent seeds. The measured
leakage factor follows.

\begin{equation}
\mathrm{LF}(B) =
\frac{k_\infty^\mathrm{assembly}(B)}{k_\mathrm{eff}^\mathrm{core}(B)}
\; L_\mathrm{ax}
\label{eq:kt-lf}
\end{equation}

\noindent where the symbols have the following meaning.

\begin{itemize}
  \item $\mathrm{LF}(B)$ is the leakage factor measured at burnup $B$
        on the design's own depleted composition, dimensionless.

  \item $k_\infty^\mathrm{assembly}(B)$ is the infinite multiplication
        factor of the reflective assembly at burnup $B$,
        dimensionless.

  \item $k_\mathrm{eff}^\mathrm{core}(B)$ is the effective
        multiplication factor of the \num{32}-assembly core built with
        the same reflector thickness, loaded uniformly with the same
        depleted composition and closed by a vacuum radial boundary,
        dimensionless.

  \item $L_\mathrm{ax}$ is the axial leakage factor of
        Section~\ref{sec:meth-lax}, equal to \num{1.0289},
        dimensionless.

  \item $B$ is the burnup, in MWd/kgHM.
\end{itemize}

Two quantities are read from the measurement. The drift is the change
of the leakage factor between the two ends of the cycle and tests the
burnup assumption. The residual is the leakage factor minus the
tabulated target and tests the composition assumption. The residual at
the end of cycle is the one that moves the reported cycle length,
through the local reactivity slope of the assembly.

\begin{equation}
\Delta\mathrm{EFPD} =
\frac{\mathrm{LF}(B_\mathrm{EOC}) - k_\mathrm{target}}
     {\left. \mathrm{d}k_\infty / \mathrm{d}B \right|_{B_\mathrm{EOC}}}
\; \frac{1000}{q}
\label{eq:kt-defpd}
\end{equation}

\noindent where the symbols have the following meaning.

\begin{itemize}
  \item $\Delta\mathrm{EFPD}$ is the correction to the reported cycle
        length, in effective full-power days.

  \item $\mathrm{LF}(B_\mathrm{EOC})$ is the leakage factor of
        Equation~\eqref{eq:kt-lf} at the end-of-cycle burnup,
        dimensionless.

  \item $k_\mathrm{target}$ is the tabulated target the optimiser used
        for that design, dimensionless.

  \item $\mathrm{d}k_\infty / \mathrm{d}B$ is the local reactivity
        slope of the assembly between the two bracketing depletion
        steps, in (MWd/kgHM)$^{-1}$, negative at the end of cycle.

  \item $q$ is the core specific power, equal to
        \SI{9.9827}{\watt\per\gram} of heavy metal.

  \item $B_\mathrm{EOC}$ is the end-of-cycle burnup, in MWd/kgHM.
\end{itemize}

Table~\ref{tab:kt-burnup} reports the measurement. Every design ended
its cycle on the reactivity crossing, so no entry is set by the burnup
cap.

\begin{table}[htbp]
  \centering
  \small
  \caption[Burnup dependence of the Route B leakage factor]{Burnup
  dependence of the Route B leakage factor on the Campaign 8 designs.}
  \label{tab:kt-burnup}
  \setlength{\tabcolsep}{4.2pt}
  \begin{tabular}{ccccccccc}
    \toprule
    Design & $e$ & $t_\mathrm{refl}$ & $B_\mathrm{EOC}$ &
      $k_\mathrm{target}$ & $\mathrm{LF}_\mathrm{BOL}$ &
      $r_\mathrm{BOL}$ & drift & $\Delta\mathrm{EFPD}$ \\
    {[--]} & [wt\%] & [cm] & [MWd/kgHM] & [--] & [--] & [pcm] & [pcm] &
      [d] \\
    \midrule
     1 & 11.03 & 3.94 & 57.8 & 1.0829 & 1.0793 & $-329$ & $+33 \pm 126$ & $+56$ \\
    59 &  9.34 & 5.66 & 50.2 & 1.0820 & 1.0796 & $-220$ & $ -4 \pm 107$ & $+43$ \\
    44 &  9.23 & 5.33 & 49.8 & 1.0822 & 1.0798 & $-219$ & $-10 \pm  90$ & $+42$ \\
    21 &  7.81 & 4.92 & 42.3 & 1.0824 & 1.0800 & $-222$ & $+149 \pm 84$ & $+11$ \\
    31 &  6.59 & 4.18 & 33.4 & 1.0828 & 1.0819 & $ -85$ & $-10 \pm  85$ & $+14$ \\
    29 &  5.90 & 4.69 & 31.0 & 1.0825 & 1.0824 & $  -9$ & $-152 \pm 99$ & $+25$ \\
    23 &  4.94 & 4.81 & 23.7 & 1.0824 & 1.0830 & $ +53$ & $-67 \pm 116$ & $ +2$ \\
    42 &  4.44 & 3.76 & 19.7 & 1.0830 & 1.0836 & $ +54$ & $-22 \pm  82$ & $ -4$ \\
    47 &  4.39 & 4.29 & 19.5 & 1.0827 & 1.0830 & $ +27$ & $ +2 \pm 107$ & $ -4$ \\
    53 &  3.45 & 4.74 & 10.8 & 1.0825 & 1.0844 & $+174$ & $-84 \pm  88$ & $-14$ \\
    \bottomrule
  \end{tabular}
\end{table}

The burnup assumption holds. No single drift is resolved at two
standard deviations, the largest ratio of value to uncertainty being
\num{1.78} on design 21. The ten drifts are ten independent
measurements of the same physical quantity, so they are combined by
inverse-variance weighting.

\begin{equation}
\overline{\Delta\mathrm{LF}} = -13 \pm \SI{30}{pcm},
\qquad
\chi^2 = 6.73 \ \text{on 9 degrees of freedom}
\label{eq:kt-comp}
\end{equation}

The weighted mean is \num{0.43} standard deviations from zero. The
reduced chi-squared of \num{0.75}, with a probability of \num{0.66},
states that the observed scatter of the ten values is accounted for by
the Monte Carlo uncertainties alone, so the data support neither a
common drift nor a design-dependent one. The bound to quote is
\SI{60}{pcm} at two standard deviations on the mean, which corresponds
to \num{7} to \SI{11}{d} at the local reactivity slopes of the table,
below \SI{0.9}{\percent} of every reported cycle. The frozen leakage
factor is therefore justified over the cycle.

The composition assumption does not hold, and this is the substantive
result. The beginning-of-life residual runs from $+174$ to
\SI{-329}{pcm} and is almost perfectly ordered by enrichment.

\begin{itemize}
  \item The Pearson correlation between $r_\mathrm{BOL}$ and the
        enrichment is $-0.975$, against $-0.234$ for the reflector
        thickness, so the effect is compositional and not geometric.

  \item The linear fit is $r_\mathrm{BOL} = 341 - 62.4\,e$ in pcm,
        with $e$ the enrichment in per cent by weight, and a
        root-mean-square residual of \SI{38}{pcm}, which is inside the
        statistics of a single solve.

  \item The residual passes through zero near \SI{5.5}{\percent} by
        weight, close to the reference composition on which the table
        was built, and grows to \SI{-329}{pcm} at the
        \SI{11.03}{\percent} of design 1.
\end{itemize}

The sign is the one a two-group argument predicts. The non-leakage
probability of a bare core is governed by the product of the geometric
buckling and the migration area, and raising the enrichment raises the
absorption cross section, shrinks the migration area, and reduces the
relative leakage. A more reactive core therefore needs a smaller
leakage correction than the reference composition of the table
supplies.

The consequence for the reported results is bounded and directional. A
negative residual means the frozen table demanded more reactivity than
the depleted core required, so Route B was conservative and the
reported cycle length is a lower bound. This is the case for the six
longest-cycle designs, with corrections of $+11$ to \SI{+56}{d}, that
is \num{0.3} to \SI{1.0}{\percent}. The sign reverses below the
reference composition, and design 53 is overstated by \SI{14}{d}, or
\SI{1.3}{\percent}. Applying the correction of
Equation~\eqref{eq:kt-defpd} to all ten designs leaves the ordering of
Table~\ref{tab:c8-front} unchanged, including the \SI{27}{d} gap
between designs 42 and 47, because the objective differences on the
front are of the order of hundreds of days.

The measurement also delivers an independent replicate of the cycle
length itself. The depletion here was run in a single shot, while the
campaign ran it in chunks of two steps, and the two used different
random-number seeds, because the seed is a hash of the design
dictionary and the two dictionaries carry different key sets. Across
the nine designs whose end of cycle falls on a \SI{4}{MWd\per\kilo
\gram} step, the cycle lengths agree to a mean of
\SI{+0.09}{\percent} with a standard deviation of
\SI{0.91}{\percent}. This is the first direct measurement of the
Monte Carlo uncertainty on cycle length in this work, and it also
confirms that the chunked restart of the campaign carries no
systematic bias.

Three limits of the measurement should be stated.

\begin{enumerate}
  \item The depleted composition is loaded uniformly into all
        \num{32} assemblies. This reproduces the construction of the
        calibration table exactly, which is what the check requires,
        but it does not represent the radial burnup distribution of an
        operating core, in which the periphery burns more slowly than
        the centre.

  \item The end-of-cycle leakage factor is interpolated linearly
        between the two bracketing depletion steps, which are up to
        \SI{4}{MWd\per\kilo\gram} apart. The interpolation error is
        second order in that interval and is not separated from the
        statistical uncertainty.

  \item Design 53 ends its cycle at \SI{10.8}{MWd\per\kilo\gram},
        inside the beginning-of-life block, so its crossing is
        interpolated across the \SI{6}{MWd\per\kilo\gram} interval
        between the depletion points at \num{7.5} and
        \SI{13.5}{MWd\per\kilo\gram}. Its cycle length is
        correspondingly more sensitive to eigenvalue noise, and it is
        the one design whose two replicates differ by more than
        \SI{1.5}{\percent}, at \SI{-4.8}{\percent}. The coarse
        beginning-of-life block is adequate for the long-cycle
        designs and marginal for the short-cycle ones.
\end{enumerate}

%% AUTHOR: whether the composition dependence of the leakage factor
%% should be reported as a limitation of Route B, as done above, or
%% closed by refitting the table as a function of enrichment as well
%% as reflector thickness in any follow-on work. The measured slope of
%% -62.4 pcm per weight per cent is the coefficient such a refit would
%% need.

%% AUTHOR: whether the conservative direction of the correction on the
%% six longest-cycle designs is worth carrying into the comparison
%% with the reported refuelling interval of
%% Section~\ref{sec:res-c8-mission}, where a lower bound on the cycle
%% length strengthens rather than weakens the argument.
